\chapter{Metabolism}

\section*{Definition and Overview}
Metabolism encompasses all the chemical reactions that occur within living organisms to maintain life. It comprises two complementary processes: catabolism (the breakdown of molecules to release energy) and anabolism (the synthesis of complex molecules from simpler precursors, requiring energy). Metabolic pathways transform nutrients (carbohydrates, fats, proteins) into adenosine triphosphate (ATP), the universal energy currency of cells, and provide building blocks for cellular structures. Dysregulation of metabolism underlies many common diseases, including diabetes, obesity, cardiovascular disease, and inborn errors of metabolism.

\section*{Epidemiology}
Metabolic diseases are among the leading causes of morbidity and mortality worldwide. Type 2 diabetes affects over 500 million people globally; obesity affects over 1 billion. In Australia, two-thirds of adults are overweight or obese, and one in twenty has diabetes. Inborn errors of metabolism, while individually rare (collective incidence ~1 in 800--1,000 births), are a significant cause of neonatal and childhood morbidity. Metabolic syndrome affects ~30\% of adults. The prevalence of metabolic diseases is rising due to ageing populations, sedentary lifestyles, and dietary changes.

\section*{Aetiology and Pathophysiology}
\begin{itemize}
    \item \textbf{Energy balance:} obesity results from chronic positive energy balance; weight loss requires negative energy balance
    \item \textbf{Macronutrient metabolism:}
    \begin{itemize}
        \item \textbf{Carbohydrates:} glycolysis, gluconeogenesis, glycogenolysis, glycogenesis, pentose phosphate pathway; regulated by insulin, glucagon, cortisol, catecholamines
        \item \textbf{Lipids:} lipolysis, $\beta$-oxidation, ketogenesis, lipogenesis, cholesterol synthesis; regulated by insulin, glucagon, thyroid hormone, PPARs
        \item \textbf{Proteins:} transamination, deamination, urea cycle; amino acids as gluconeogenic or ketogenic precursors
    \end{itemize}
    \item \textbf{Key regulatory organs:}
    \begin{itemize}
        \item \textbf{Liver:} central metabolic hub; glucose production/storage, lipid synthesis/oxidation, ketogenesis, urea cycle, lipoprotein assembly
        \item \textbf{Pancreas:} insulin (anabolic) and glucagon (catabolic) from $\beta$- and $\alpha$-cells
        \item \textbf{Adipose tissue:} energy storage, endocrine organ (leptin, adiponectin, resistin), lipid buffering
        \item \textbf{Muscle:} major site of glucose disposal and protein metabolism
        \item \textbf{Brain:} obligate glucose consumer (except during prolonged starvation/ketosis); hypothalamic appetite regulation
    \end{itemize}
    \item \textbf{Mitochondrial function:} oxidative phosphorylation, ATP synthesis, ROS production; mitochondrial dysfunction implicated in ageing, neurodegeneration, insulin resistance
    \item \textbf{Hormonal regulation:} insulin, glucagon, catecholamines, cortisol, thyroid hormone, growth hormone, sex steroids, FGF21, GDF15
    \item \textbf{Circadian regulation:} metabolic rhythms synchronised to feeding-fasting cycles; disruption promotes metabolic disease
    \item \textbf{Gut microbiome:} ferments fibre to SCFAs; influences energy harvest, inflammation, bile acid metabolism, GLP-1 secretion
\end{itemize}

\section*{Clinical Features}
Metabolic disease presentations vary widely:
\begin{itemize}
    \item \textbf{Diabetes/hyperglycaemia:} polyuria, polydipsia, weight loss, fatigue, blurred vision, recurrent infections, ketoacidosis (type 1), hyperosmolar state (type 2)
    \item \textbf{Hypoglycaemia:} adrenergic symptoms (tremor, sweating, tachycardia), neuroglycopenic symptoms (confusion, seizures, coma)
    \item \textbf{Obesity:} increased adiposity, breathlessness, joint pain, sleep apnoea, stigma
    \item \textbf{Dyslipidaemia:} usually asymptomatic; xanthomas, xanthelasmas, arcus corneae in severe cases; pancreatitis with severe hypertriglyceridaemia
    \item \textbf{Inborn errors of metabolism (acute):} neonatal encephalopathy, vomiting, lethargy, seizures, hypoglycaemia, hyperammonemia, acidosis, unusual odours
    \item \textbf{Inborn errors (chronic):} developmental delay, failure to thrive, organomegaly, dysmorphism, movement disorders, regression
    \item \textbf{Thyroid dysfunction:} hyperthyroidism (weight loss, tachycardia, tremor, heat intolerance) / hypothyroidism (weight gain, fatigue, cold intolerance, bradycardia)
    \item \textbf{Adrenal disorders:} Cushing (central obesity, hypertension, glucose intolerance) / Addison (hypotension, hyponatraemia, hyperkalaemia, hypoglycaemia)
\end{itemize}

\section*{Differential Diagnosis}
\begin{itemize}
    \item \textbf{Type 1 vs type 2 diabetes:} age, BMI, ketosis, autoantibodies (GAD, IA-2, ZnT8, insulin), C-peptide
    \item \textbf{Secondary diabetes:} pancreatic disease, Cushing, acromegaly, haemochromatosis, drug-induced (steroids, antipsychotics), monogenic (MODY)
    \item \textbf{Obesity:} endocrine causes (hypothyroidism, Cushing, hypothalamic), genetic syndromes (Prader-Willi, Bardet-Biedl, MC4R), medications
    \item \textbf{Hyperlipidaemia:} primary (familial hypercholesterolaemia, familial combined, familial hypertriglyceridaemia) vs secondary (diabetes, hypothyroidism, nephrotic syndrome, alcohol, medications)
    \item \textbf{Metabolic acidosis:} diabetic ketoacidosis, lactic acidosis, renal tubular acidosis, toxic alcohols, inborn errors (organic acidaemias)
    \item \textbf{Hyperammonaemia:} urea cycle disorders, organic acidaemias, fatty acid oxidation defects, liver failure, medications (valproate)
\end{itemize}

\section*{When to Seek Medical Care}
\begin{itemize}
    \item Symptoms of diabetes, hypoglycaemia, or thyroid dysfunction
    \item Unexplained weight loss or gain, fatigue, polyuria
    \item Family history of early cardiovascular disease, diabetes, or sudden death
    \item Neonate/infant with poor feeding, vomiting, lethargy, seizures, developmental regression -- urgent metabolic screen
    \item Acute metabolic decompensation in known metabolic disorder (intercurrent illness, fasting)
\end{itemize}

\textbf{Call 000 for:}
\begin{itemize}
    \item Severe hypoglycaemia with altered consciousness or seizures
    \item Diabetic ketoacidosis (Kussmaul breathing, dehydration, altered consciousness)
    \item Hyperosmolar hyperglycaemic state (severe dehydration, hypernatraemia, altered consciousness)
    \item Adrenal crisis (hypotension, hyponatraemia, hyperkalaemia, hypoglycaemia)
    \item Myxoedema coma (hypothermia, bradycardia, altered consciousness)
    \item Acute metabolic crisis in infant/child (hyperammonaemia, severe acidosis)
\end{itemize}

\section*{Diagnosis and Investigations}
\begin{itemize}
    \item \textbf{Basic metabolic panel:} glucose, HbA1c, electrolytes, renal function, calcium, phosphate, magnesium
    \item \textbf{Lipid profile:} total cholesterol, LDL, HDL, triglycerides; apoB, non-HDL, Lp(a) for risk stratification
    \item \textbf{Hormonal:} TSH, free T4, free T3, cortisol (AM), ACTH, 17-OHP, androgens, SHBG, IGF-1
    \item \textbf{Diabetes-specific:} autoantibodies, C-peptide, ketones, urinary ACR, eGFR, eye screen, foot exam
    \item \textbf{Inborn errors (acute):} blood gas, ammonia, lactate, glucose, ketones, acylcarnitine profile, amino acids, organic acids, urine reducing substances
    \item \textbf{Genetic testing:} gene panels for monogenic diabetes, familial hypercholesterolaemia, inborn errors; whole exome/genome for undiagnosed metabolic disease
    \item \textbf{Imaging:} liver ultrasound (NAFLD), DEXA (body composition), CT/MRI (pituitary, adrenal, pancreas)
    \item \textbf{Functional tests:} OGTT, insulin tolerance test, glucagon stimulation, water deprivation test, exercise testing
\end{itemize}

\section*{Management}
\begin{itemize}
    \item \textbf{Lifestyle (foundation for acquired metabolic disease):} Mediterranean/DASH diet, regular physical activity (aerobic + resistance), weight management, sleep hygiene, stress reduction, smoking cessation, alcohol moderation
    \item \textbf{Diabetes:} metformin first-line; GLP-1 RAs and SGLT2 inhibitors for cardiorenal benefit and weight loss; insulin for type 1 and advanced type 2; technology (CGM, pumps, automated insulin delivery)
    \item \textbf{Obesity:} intensive behavioural therapy; pharmacotherapy (GLP-1 RAs, GIP/GLP-1 dual agonists, naltrexone/bupropion); bariatric surgery (sleeve, bypass) for BMI $\ge$40 or $\ge$35 with comorbidity
    \item \textbf{Dyslipidaemia:} statins (intensity by risk); ezetimibe, PCSK9 inhibitors, inclisiran for LDL not at goal; fibrates/omega-3 for hypertriglyceridaemia
    \item \textbf{Thyroid:} levothyroxine for hypothyroidism; antithyroid drugs, radioiodine, or surgery for hyperthyroidism
    \item \textbf{Adrenal:} glucocorticoid +/- mineralocorticoid replacement for Addison; surgery for Cushing/phaeochromocytoma
    \item \textbf{Inborn errors:} substrate restriction, cofactor supplementation, enzyme replacement, chaperone therapy, substrate reduction, gene therapy, liver transplant; emergency protocols for intercurrent illness
    \item \textbf{Multidisciplinary care:} metabolic physician, endocrinologist, dietitian, genetic counsellor, nurse specialist, psychologist
\end{itemize}

\section*{Complications}
\begin{itemize}
    \item \textbf{Diabetes:} microvascular (retinopathy, nephropathy, neuropathy), macrovascular (CVD, stroke, PAD), foot ulceration/amputation, hypoglycaemia, infections
    \item \textbf{Obesity:} CVD, type 2 diabetes, NAFLD/MAFLD, OSA, osteoarthritis, reproductive dysfunction, cancer, depression
    \item \textbf{Dyslipidaemia:} premature atherosclerosis, coronary artery disease, stroke, pancreatitis (severe hypertriglyceridaemia)
    \item \textbf{Inborn errors:} intellectual disability, organ failure, movement disorders, sudden death, crisis-related brain injury
    \item \textbf{Thyroid/adrenal:} cardiovascular, bone, reproductive, neuropsychiatric sequelae
\end{itemize}

\section*{Prognosis and Outcomes}
Acquired metabolic diseases are largely preventable and manageable. Intensive lifestyle intervention can prevent or delay type 2 diabetes (58\% reduction in DPP), induce remission of early type 2 diabetes (DiRECT trial), and reverse NAFLD. Modern pharmacotherapy (GLP-1 RAs, SGLT2 inhibitors) provides cardiorenal protection beyond glucose lowering. For inborn errors, newborn screening enables presymptomatic diagnosis and treatment, dramatically improving outcomes; novel therapies (enzyme replacement, gene therapy, mRNA) are transforming previously fatal disorders. Lifelong adherence and monitoring are essential.

\section*{Prevention and Self-Care}
\begin{itemize}
    \item Maintain healthy weight through balanced diet and regular activity from childhood
    \item Follow Australian Dietary Guidelines / Mediterranean pattern
    \item 150--300 min/week moderate aerobic activity + 2 sessions resistance training
    \item Limit ultra-processed foods, added sugars, saturated fats, alcohol
    \item Don't smoke; avoid second-hand smoke
    \item Regular health checks: BP, lipids, glucose, waist from age 40 (or earlier if risk factors)
    \item Know family history of diabetes, CVD, lipid disorders, sudden death
    \item Newborn screening for inborn errors (heel prick test) -- ensure completed
    \item If diagnosed with metabolic condition: adhere to treatment, monitor as advised, carry emergency information (medical alert for diabetes, adrenal insufficiency, inborn errors)
    \item Sick-day rules for diabetes and inborn errors: never stop insulin; increase monitoring; seek care early
\end{itemize}

\section*{Sources and Further Reading}
The following reputable sources provide further information for healthcare students and patients seeking reliable, current advice about metabolism and metabolic disorders.

\begin{itemize}
    \item Australian Diabetes Society. \textit{Clinical guidelines.} https://diabetessociety.com.au/
    \item Endocrine Society of Australia. \textit{Position statements.} https://www.endocrinesociety.org.au/
    \item Human Genetics Society of Australasia. \textit{Inborn errors of metabolism resources.} https://www.hgsa.org.au/
    \item Metabolic Dietary Disorders Association (MDDA). \textit{Support for inborn errors.} https://www.mdda.org.au/
    \item Royal Australasian College of Physicians. \textit{Obesity management guidelines.} https://www.racp.edu.au/
    \item Baker Heart and Diabetes Institute. \textit{Metabolic research and resources.} https://www.baker.edu.au/
    \item Healthdirect Australia. \textit{Metabolism and metabolic disorders.} https://www.healthdirect.gov.au/
    \item National Organization for Rare Disorders (NORD). \textit{Inborn errors of metabolism.} https://rarediseases.org/
    \item Online Mendelian Inheritance in Man (OMIM). \textit{Genetic metabolic disorders.} https://www.omim.org/
\end{itemize}